\section*{第 \ref{ch:g12:glucose-and-diabetes} 章——血糖与糖尿病}

\begin{solution}{exo:g12:glucose-and-diabetes:1}
约 \qty{1}{g/L}，即 \qty{5.5}{mmol/L}。低于 \qty{0.5}{g/L} 时，
脑功能衰竭；高于 \qty{1.8}{g/L} 时，葡萄糖进入尿液（
长期高于 \qty{1.3}{g/L} 会损伤血管）。
\end{solution}

\begin{solution}{exo:g12:glucose-and-diabetes:2}
\omterm{prop:g12:glucose-and-diabetes:hormones}{胰岛素}由 $\beta$ \omterm{def:g10:cells-common-unit:cell}{细胞}分泌，可降低\omterm{def:g12:glucose-and-diabetes:glycaemia}{血糖}（经肌肉、脂肪和肝脏
摄取并储存葡萄糖）；\omterm{prop:g12:glucose-and-diabetes:hormones}{胰高血糖素}由 $\alpha$ \omterm{def:g10:cells-common-unit:cell}{细胞}分泌，
可升高\omterm{def:g12:glucose-and-diabetes:glycaemia}{血糖}（经肝脏释放葡萄糖）。
\end{solution}

\begin{solution}{exo:g12:glucose-and-diabetes:3}
肝脏以\omterm{def:g10:exercise-and-energy:glycogen}{糖原}的形式储存（它还能由
\omterm{def:g10:chemistry-of-life:families}{氨基酸}和乳酸生成新的葡萄糖）。
\end{solution}

\begin{solution}{exo:g12:glucose-and-diabetes:4}
1 型：$\beta$ \omterm{def:g10:cells-common-unit:cell}{细胞}被破坏，\omterm{prop:g12:glucose-and-diabetes:hormones}{胰岛素}缺失。2 型：
\omterm{prop:g12:glucose-and-diabetes:hormones}{胰岛素}仍存在，但组织对其反应不佳，
$\beta$ \omterm{def:g10:cells-common-unit:cell}{细胞}最终也会衰竭。
\end{solution}

\begin{solution}{exo:g12:glucose-and-diabetes:5}
健康人约为 \qty{0.95}{g/L}；糖耐量受损约为 \qty{1.7}{g/L}
（介于 1.4 与 2 之间）；2 型约为 \qty{2.7}{g/L}（高于 2）。
\end{solution}

\begin{solution}{exo:g12:glucose-and-diabetes:6}
肝脏从\omterm{def:g10:exercise-and-energy:glycogen}{糖原}中释放葡萄糖，释放速率与脑的消耗速率一致，
并受\omterm{prop:g12:glucose-and-diabetes:hormones}{胰高血糖素}调控，因此葡萄糖库维持在
\qty{5}{g} 左右，尽管其内容每小时都会更新。
\end{solution}

\begin{solution}{exo:g12:glucose-and-diabetes:7}
导管把消化酶输送到肠道；堵塞导管后，
\omterm{def:g12:glucose-and-diabetes:glycaemia}{血糖}仍然正常，说明胰腺对\omterm{def:g12:glucose-and-diabetes:glycaemia}{血糖}的作用并非通过
消化实现的。切除整个器官即去除了\omterm{prop:g12:glucose-and-diabetes:hormones}{胰岛}：\omterm{prop:g12:glucose-and-diabetes:hormones}{胰岛素}
在胰腺中合成，经血液而非
导管运送至靶器官。
\end{solution}

\begin{solution}{exo:g12:glucose-and-diabetes:8}
\omterm{prop:g12:glucose-and-diabetes:hormones}{胰高血糖素}通过促使肝脏分解\omterm{def:g10:exercise-and-energy:glycogen}{糖原}发挥作用；
一夜禁食后，\omterm{def:g10:exercise-and-energy:glycogen}{糖原}储备仍在，\omterm{def:g12:glucose-and-diabetes:glycaemia}{血糖}上升；禁食三天后，
\omterm{def:g10:exercise-and-energy:glycogen}{糖原}已耗尽，\omterm{prop:g12:glucose-and-diabetes:hormones}{胰高血糖素}便无物可释放。
\end{solution}

\begin{solution}{exo:g12:glucose-and-diabetes:9}
\omterm{def:g12:glucose-and-diabetes:glycaemia}{血糖}高于约 \qty{1.8}{g/L} 时，肾脏无法重吸收全部滤过的葡萄糖，
葡萄糖随尿液排出，并通过渗透作用带走水分。
失去的水分须得到补充，于是产生口渴。
\end{solution}

\begin{solution}{exo:g12:glucose-and-diabetes:10}
\omterm{prop:g12:glucose-and-diabetes:hormones}{胰岛素}使葡萄糖进入肌肉和脂肪\omterm{def:g10:cells-common-unit:cell}{细胞}，并停止
肝脏释放葡萄糖，但没有新的葡萄糖进入血液，\omterm{def:g12:glucose-and-diabetes:glycaemia}{血糖}
在一两小时内降至 \qty{0.6}{g/L} 以下，并伴有出汗、颤抖
和意识模糊；此时立即摄入糖分即可纠正。
\end{solution}

\begin{solution}{exo:g12:glucose-and-diabetes:11}
运动中的肌肉通过一条不需要\omterm{prop:g12:glucose-and-diabetes:hormones}{胰岛素}的途径从血液摄取葡萄糖，
经过训练的肌肉在两次运动之间能恢复对\omterm{prop:g12:glucose-and-diabetes:hormones}{胰岛素}的敏感性；
运动期间及运动后\omterm{def:g12:glucose-and-diabetes:glycaemia}{血糖}下降，
组织的\omterm{prop:g12:glucose-and-diabetes:hormones}{胰岛素}抵抗（该病的病因）得以减轻。
\end{solution}

\begin{solution}{exo:g12:glucose-and-diabetes:12}
无\omterm{prop:g12:glucose-and-diabetes:hormones}{胰岛素}：为 1 型，\omterm{prop:g12:glucose-and-diabetes:hormones}{胰岛}的 $\beta$ \omterm{def:g10:cells-common-unit:cell}{细胞}被破坏，不再分泌
任何\omterm{prop:g12:glucose-and-diabetes:hormones}{胰岛素}。\omterm{prop:g12:glucose-and-diabetes:hormones}{胰岛素}水平为正常的三倍：为 2 型，\omterm{prop:g12:glucose-and-diabetes:hormones}{胰岛}持续
满负荷分泌，而靶组织已不再对其做出反应。
\end{solution}

\begin{solution}{exo:g12:glucose-and-diabetes:13}
健康人：\omterm{def:g12:glucose-and-diabetes:glycaemia}{血糖}在 30--60 分钟时升至约 \qty{1.3}{g/L}，2 小时内
恢复正常；\omterm{prop:g12:glucose-and-diabetes:hormones}{胰岛素}随之升高、随之下降。1 型未治疗者：\omterm{def:g12:glucose-and-diabetes:glycaemia}{血糖}
升至 \qty{2}{g/L} 以上并持续不降；\omterm{prop:g12:glucose-and-diabetes:hormones}{胰岛素}水平始终为零——
无物清除葡萄糖。2 型：\omterm{def:g12:glucose-and-diabetes:glycaemia}{血糖}升得很高，下降缓慢；\omterm{prop:g12:glucose-and-diabetes:hormones}{胰岛素}
水平高于正常并持续处于高位——
\omterm{def:g11:hormones-and-reproduction:hormone}{激素}量大，作用却微弱。
\end{solution}

\begin{solution}{exo:g12:glucose-and-diabetes:14}
$5 + 80 = \qty{85}{g}$ 分布于 \qty{5}{L} 中，即 \qty{17}{g/L}。这种情况从未发生，
因为吸收需要一两个小时，随着葡萄糖进入体内，
\omterm{prop:g12:glucose-and-diabetes:hormones}{胰岛素}将其转入肝\omterm{def:g10:exercise-and-energy:glycogen}{糖原}、肌\omterm{def:g10:exercise-and-energy:glycogen}{糖原}和脂肪，同时
脑和其他组织也会消耗一部分，因此葡萄糖库
中的量从不会比正常水平多出几克。
\end{solution}

\begin{solution}{exo:g12:glucose-and-diabetes:15}
两者均具备：感受器（\omterm{prop:g12:glucose-and-diabetes:hormones}{胰岛}\omterm{def:g10:cells-common-unit:cell}{细胞}；下丘脑和垂体）、设定点、
效应器，以及移除信号的负\omterm{def:g11:hormones-and-reproduction:hormone}{反馈}。\omterm{def:g12:glucose-and-diabetes:glycaemia}{血糖}
由两种作用相反的\omterm{def:g11:hormones-and-reproduction:hormone}{激素}同时作用于同一器官，因此该变量
可被主动地向上或向下推动；睾酮的调节环路只有一种
\omterm{def:g11:hormones-and-reproduction:hormone}{激素}，主要通过增减其分泌量来纠正偏差。两种作用相反的
效应器能对双向快速变化的变量实现更快、更精确的
调节。
\end{solution}

\begin{solution}{pb:g12:glucose-and-diabetes:1}
\textbf{1.} \qty{5}{g}；$\qty{1}{g/L} \div \qty{180}{g/mol} = \qty{5.55}{mmol/L}$。

\textbf{2.} 一小时。

\textbf{3.} $100/5 = 20$ 小时。此后，肝脏
由\omterm{def:g10:chemistry-of-life:families}{氨基酸}（及乳酸）生成新的葡萄糖，代价是消耗机体的\omterm{def:g11:gene-expression:protein}{蛋白质}。

\textbf{4.} 肌肉\omterm{def:g10:cells-common-unit:cell}{细胞}缺乏能将\omterm{def:g10:exercise-and-energy:glycogen}{糖原}释放为游离葡萄糖、使其进入血液的酶；
其\omterm{def:g10:exercise-and-energy:glycogen}{糖原}只能就地燃烧，
供自身收缩之用。

\textbf{5.} $5 + 60 = \qty{65}{g}$，溶于 \qty{5}{L} 中：\qty{13}{g/L}。
\omterm{prop:g12:glucose-and-diabetes:hormones}{胰岛素}促使葡萄糖转化为肝脏和肌肉中的\omterm{def:g10:exercise-and-energy:glycogen}{糖原}，并转入脂肪，
部分葡萄糖被组织氧化分解；\omterm{def:g12:glucose-and-diabetes:glycaemia}{血糖}峰值约为 \qty{1.4}{g/L}。

\textbf{6.} 空腹时 \qty{1.15}{g/L}（低于 1.26），而
120 分钟时 \qty{1.85}{g/L}（介于 1.4 与 2 之间）：糖耐量受损，
处于\omterm{def:g12:glucose-and-diabetes:diabetes}{糖尿病}的临界状态。

\textbf{7.} 否：\omterm{prop:g12:glucose-and-diabetes:hormones}{胰岛素}水平已达健康峰值的两倍，组织对它的
反应已失效——即\omterm{prop:g12:glucose-and-diabetes:hormones}{胰岛素}抵抗。

\textbf{8.} 2 型（或其预警阶段）。预期人群：成年人、超重、
缺乏运动，且常有\omterm{def:g12:glucose-and-diabetes:diabetes}{糖尿病}患病亲属。

\textbf{9.} \omterm{def:g12:glucose-and-diabetes:glycaemia}{血糖}起始高于 \qty{1.5}{g/L}，并升至
高于 \num{2.5} 且不下降；全程\omterm{prop:g12:glucose-and-diabetes:hormones}{胰岛素}接近于零。

\textbf{10.} 减重和规律运动：二者可恢复
组织对\omterm{prop:g12:glucose-and-diabetes:hormones}{胰岛素}的敏感性，并使运动中的肌肉
无需\omterm{prop:g12:glucose-and-diabetes:hormones}{胰岛素}即可摄取葡萄糖。

\textbf{11.} $90/10 = 9$ 单位。

\textbf{12.} 该剂量可处理 \qty{90}{g} 葡萄糖，但只有 \qty{45}{g}
到达血液：多余\omterm{prop:g12:glucose-and-diabetes:hormones}{胰岛素}使\omterm{def:g12:glucose-and-diabetes:glycaemia}{血糖}下降相当于
血池与储备中 \qty{45}{g} 的量——\omterm{def:g12:glucose-and-diabetes:glycaemia}{血糖}远低于
\qty{0.6}{g/L}：出汗、颤抖、意识模糊；须立即补充糖分。

\textbf{13.} 若无任何清除，则 $(5 + 90)/5 = \qty{19}{g/L}$；实际
上仅为数 \unit{g/L}。高于 \qty{1.8}{g/L} 时，肾脏使
葡萄糖随水排入尿液：尿量增多，口渴。

\textbf{14.} 两餐之间，若无调节，肝脏将无约束地释放葡萄糖
（\omterm{prop:g12:glucose-and-diabetes:hormones}{胰高血糖素}起作用，且无物拮抗），组织则会
大量燃烧脂肪并产生酸性产物：需要一定基础水平的\omterm{prop:g12:glucose-and-diabetes:hormones}{胰岛素}，以维持
基本平衡，而不仅是为了进餐。

\textbf{15.} \omterm{def:g12:glucose-and-diabetes:glycaemia}{血糖}仪测量\omterm{def:g12:glucose-and-diabetes:glycaemia}{血糖}的方式与 $\beta$ \omterm{def:g10:cells-common-unit:cell}{细胞}一致；
\omterm{prop:g12:glucose-and-diabetes:hormones}{胰岛素}笔输送的是其本应分泌的\omterm{prop:g12:glucose-and-diabetes:hormones}{胰岛素}。二者均无法
替代持续、逐分钟的调节：患者
每天只测量数次并估算，而\omterm{def:g10:cells-common-unit:cell}{细胞}则是持续测量、
并精确响应。

\textbf{16.} 2 型，其 $\beta$ \omterm{def:g10:cells-common-unit:cell}{细胞}尚存且可被刺激。在
1 型中则无 $\beta$ \omterm{def:g10:cells-common-unit:cell}{细胞}可供刺激。

\textbf{17.} 组织对\omterm{prop:g12:glucose-and-diabetes:hormones}{胰岛素}的反应——即调节回路中的效应器一侧：
脂肪减少、肌肉活跃后，正常的\omterm{prop:g12:glucose-and-diabetes:hormones}{胰岛素}信号
又能引起正常的摄取，$\beta$ \omterm{def:g10:cells-common-unit:cell}{细胞}由此得以减负。

\textbf{18.} 二者各自响应同一变量的相反偏差而分泌：
\omterm{def:g12:glucose-and-diabetes:glycaemia}{血糖}升高时促使\omterm{prop:g12:glucose-and-diabetes:hormones}{胰岛素}分泌、抑制\omterm{prop:g12:glucose-and-diabetes:hormones}{胰高血糖素}分泌，
\omterm{def:g12:glucose-and-diabetes:glycaemia}{血糖}降低时则相反。在 \qty{1}{g/L} 时，\omterm{prop:g12:glucose-and-diabetes:hormones}{胰岛}
两种\omterm{def:g11:hormones-and-reproduction:hormone}{激素}均以低而平衡的速率分泌——即其静息状态。

\textbf{19.} \omterm{def:g12:glucose-and-diabetes:glycaemia}{血糖}过低数分钟内即可致使脑部受损，过高则
经多年缓慢损伤：急性危险有多重相互独立的
防护，而慢性危险仅有一重防护——这正是\omterm{prop:g12:glucose-and-diabetes:hormones}{胰岛素}功能
单独失常即足以引发疾病，也是\omterm{def:g12:glucose-and-diabetes:diabetes}{糖尿病}
如此常见的原因。

\textbf{20.} 血液中约含 \qty{5}{g} 葡萄糖；肝脏的
\omterm{def:g10:exercise-and-energy:glycogen}{糖原}可维持约 20 小时；第二部分中的患者糖耐量受损，
处于 2 型\omterm{def:g12:glucose-and-diabetes:diabetes}{糖尿病}的边缘；第三部分的剂量为 9
单位。
\end{solution}
